\section{A Brownian last-passage identity}

Write $\Phi$ for the standard normal distribution function and
$\PhiBar=1-\Phi$.  Throughout, $\GammaLaw(\alpha,\theta)$ denotes the
gamma distribution with shape $\alpha$ and scale $\theta$.

\begin{lemma}[Exact last-passage law]\label{lem:last-passage}
Let $B$ be standard Brownian motion and let $a>0$.  Define
\begin{align}
  L_a^{>} &= \sup\{s\ge0:B(s)-as>0\},\label{eq:L-strict}\\
  L_a^{\ge} &= \sup\{s\ge0:B(s)-as\ge0\}.
  \label{eq:L-nonstrict}
\end{align}
Both variables are finite almost surely and have the same distribution.  For
$t\ge0$,
\begin{equation}\label{eq:last-passage-cdf}
  \P(L_a^{>}\le t)
  =\P(L_a^{\ge}\le t)
  =2\Phi(a\sqrt t)-1.
\end{equation}
Equivalently, if $G\sim N(0,1)$, then
\begin{equation}\label{eq:last-passage-chi-square}
  L_a^{>}\ \stackrel{\mathrm d}{=}\
  L_a^{\ge}\ \stackrel{\mathrm d}{=}\frac{G^2}{a^2}.
\end{equation}
\end{lemma}

\begin{proof}
Set $R(s)=B(s)-as$.  Since $B(s)/s\to0$ almost surely, $R(s)<0$ for all
sufficiently large $s$, so both last-passage times are finite.

We first recall the hitting probability from a negative starting point.  If
$R(0)=y<0$ and $\tau_0=\inf\{s\ge0:R(s)=0\}$, then
\begin{equation}\label{eq:hitting-probability}
  \P_y(\tau_0<\infty)=e^{2ay}.
\end{equation}
Indeed, $e^{2aR(s)}$ is a local martingale.  Stopping it on first hitting
$0$ or $-m$, and then letting $m\to\infty$, gives
\eqref{eq:hitting-probability}.  This is the standard one-sided hitting
formula for Brownian motion with drift; see, for example,
\cite[Chapter~3]{KaratzasShreve1991}.

Fix $t>0$.  Conditional on $R(t)=y<0$, a future visit to the nonnegative
half-line has probability $e^{2ay}$.  If $y>0$, a future positive value occurs
with probability one by continuity.  Starting from $y=0$, the positive
half-line is regular for Brownian motion with drift, so a future positive value
again occurs with probability one.  The distinction at $y=0$ is immaterial in
any case because $R(t)$ has a continuous distribution.  The same calculation
therefore applies to the strict and non-strict last-passage times:
\begin{equation}\label{eq:last-passage-tail-decomp}
  \P(L_a^{\diamond}>t)
  =\P(R(t)\ge0)
   +\E\!\left[e^{2aR(t)};R(t)<0\right],
  \qquad \diamond\in\{>,\ge\}.
\end{equation}
Now $R(t)\sim N(-at,t)$, so
\begin{equation}\label{eq:first-half-tail}
  \P(R(t)\ge0)=\Phi(-a\sqrt t).
\end{equation}
Completing the square in the second term gives
\begin{align}
  \E\!\left[e^{2aR(t)};R(t)<0\right]
  & =\int_{-\infty}^{0}
  \frac{1}{\sqrt{2\pi t}}
  \exp\!\left(2ay-\frac{(y+at)^2}{2t}\right)\dd y\notag\\
  & =\int_{-\infty}^{0}
  \frac{1}{\sqrt{2\pi t}}
  \exp\!\left(-\frac{(y-at)^2}{2t}\right)\dd y\notag\\
  & =\Phi(-a\sqrt t).
  \label{eq:second-half-tail}
\end{align}
Thus
\begin{equation}\label{eq:last-passage-tail}
  \P(L_a^{\diamond}>t)=2\Phi(-a\sqrt t),
\end{equation}
which is equivalent to \eqref{eq:last-passage-cdf}.  Finally,
\[
  \P(G^2/a^2\le t)=\P(|G|\le a\sqrt t)=2\Phi(a\sqrt t)-1,
\]
proving \eqref{eq:last-passage-chi-square}.
\end{proof}

\begin{remark}[Time inversion]\label{rem:time-inversion}
Lemma~\ref{lem:last-passage} also follows immediately from Brownian time
inversion and the reflection principle.  If
$\widetilde B(u)=uB(1/u)$, then $\widetilde B$ is Brownian motion, and
\[
 \{L_a^{>}\le t\}
 =\left\{\sup_{0\le u\le1/t}\widetilde B(u)\le a\right\}.
\]
The reflection principle then gives
$2\Phi(a\sqrt t)-1$.  See \cite{RevuzYor1999,BorodinSalminen2002} for
classical formulations of these identities.
\end{remark}
